\documentclass[10pt,a4paper]{amsart}
\usepackage[margin=19mm]{geometry}
\usepackage{amsmath,amssymb,mathtools}
\usepackage[scr=rsfs,cal=euler]{mathalfa}
\usepackage{xfrac}
\usepackage[unicode,hidelinks,bookmarks=false]{hyperref}
\newcommand{\ie}{i.e.\ }
\newcommand{\cf}{cf.\ }
\newcommand{\resp}{respectively\ }
\newcommand{\Subsection}{\subsection}
\providecommand{\nobreakdash}{\nobreak\hbox{-}}
\setlength{\emergencystretch}{2em}
\multlinegap0pt
\usepackage{amssymb}
\usepackage{multirow}
\usepackage{booktabs}
\usepackage{caption}
\usepackage{xcolor}
\usepackage{float}
\newtheorem{lemma}{Lemma}
\title{Sinh regularized Lagrangian nonuniform sampling series}
\author{Haixin Jiang, Xinyu Chen, Liang Chen}
\date{Source: arXiv:2601.17685v2 (CC BY 4.0)}
\begin{document}
\maketitle

\noindent\emph{Source and licence.} Sinh regularized Lagrangian nonuniform sampling series. Version arXiv:2601.17685v2 (CC BY 4.0), distributed by arXiv under the Creative Commons Attribution 4.0 International licence, http://creativecommons.org/licenses/by/4.0/, as recorded in the arXiv licence metadata for this exact version and shown on the abstract page https://arxiv.org/abs/2601.17685v2. Full licence notice: \texttt{licenses/arxiv-2601.17685-CC-BY-4.0.txt}.\par\smallskip

\noindent\textit{Editorial scope.} Counted unit: the article's lemma L11 (source0.tex lines 164--172). The article's complete original proof of it is delivered with it (source0.tex lines 174--191, one \texttt{proof} environment of 18 lines). No mathematics is written, repaired or re-derived: the statement and the proof are the article's own lines, in the article's own notation, and no retained line is substituted or re-flowed.  No further source-local context is needed: every cross-reference of every retained line resolves inside the two delivered spans.  Nothing was omitted from any retained span, so the package carries no omission record.  No retained line contains a citation, so the wrapper carries no bibliography and no bibliography entry.  No retained line contains a figure environment, an image-inclusion command or drawing code, so no image is delivered and the package declares no asset dependency.  Editorial formatting only, in addition: an AMS \texttt{amsart} preamble that restates the article's own declarations and macros used by the retained lines, namely the environments \texttt{lemma} on separate counters, together with the standard packages those lines need.  The retained environments print in this document as Lemma 1 in order of appearance, whereas the article prints the same items with the numbering of its own full document.  The complete native source of the article as distributed in the arXiv e-print for this exact version is bundled unchanged as \texttt{sources/193338/source0.tex}.
\begin{lemma}\label{L11}
Let $ N\in \mathbb{N}$, $0<\delta<\pi$, $f\in \mathcal{B}_{\delta}(\mathbb{R})$ and $g_{N}\in C(\mathbb{R})\cap L^{2}(\mathbb{R})$ with $g_{N}(0)=1$ is a regularization function depending on $N$.
Assume $\{\phi_{j}\}_{j=1}^{\infty}$ is a family of basis functions in $\mathcal{B}_{\pi}(\mathbb{R})$ and $\{x_{j}\}_{j=1}^{\infty}\subseteq\mathbb{R}$ is a sample set such that for any $h\in \mathcal{B}_{\pi}(\mathbb{R})$,  the following  interpolation formula holds \begin{equation}\label{eq1}h(x)=\sum_{j=1}^{\infty}h(x_{j})\phi_{j}(x)  \quad\text{for}\quad x\in\mathbb{R}.\end{equation}  
Then, 
$$\sup_{x\in [-1,1]}|f(x) - S_{f,N}(x)|\le E_{1,N}+E_{2,N},$$
where \begin{equation}\label{eqsfn}S_{f,N}(x):=\sum_{j=1}^{N}f(x_{j})\phi_{j}(x)g_{N}(x-x_{j}),\end{equation}

$$E_{1,N}:=\sup_{x\in [-1,1]}\left\{ \frac{2}{\sqrt{2\pi}}\sum_{j=1}^{\infty}|f(x_j)\phi_{j}(x)|\left(\int_{\pi-\delta}^{\infty}|\widehat{g_{N}}(w)|dw+ \int_{-\infty}^{\delta-\pi}|\widehat{g_{N}}(w)|dw\right)\right\}$$ and $$E_{2,N}:=\sup_{x\in [-1,1]}\left\{ \sum_{j=N+1}^{\infty}|f(x_j)\phi_{j}(x)g_{N}(x-x_{j})| \right\}.$$
\end{lemma}

\begin{proof}
For $x\in \mathbb{R}$, let  $$S_{f,\infty}(x):=\sum_{j=1}^{\infty}f(x_{j})\phi_{j}(x)g_{N}(x-x_{j}),$$
and 
$$f(x) - S_{f,N}(x)=:E_{1}(x)+E_{2}(x),$$
where $$E_{1}(x):= f(x)-S_{f,\infty}(x),$$ and $$E_{2}(x):= S_{f,\infty}(x)-S_{f,N}(x).$$
It is obvious that
$\sup_{x\in[-1,1]}|E_{2}(x)|\le E_{2,N}$.
Let $$G_{N}(x)=\frac{1}{\sqrt{2\pi}}\int_{\delta-\pi}^{\pi-\delta}\widehat{g_N}(w)e^{iwx}dw$$ and $$G_{t,N}(x)=G_{N}(t-x).$$ Since ${\rm supp} \widehat{G_{t,N}} \subseteq [\delta-\pi,\pi-\delta] $, $${\rm supp}\widehat{G_{t,N}}*\widehat{f} \subseteq {\rm supp}\widehat{G_{t,N}}+ {\rm supp}\widehat{f} \subseteq[-\pi,\pi].$$ Note that $\widehat{G_{t,N}f}=\widehat{G_{t,N}}*\widehat{f} $,  we have $f(x)G_{N}(t-x)\in  \mathcal{B}_{\pi}(\mathbb{R})$. Using Eq.(\ref{eq1}), we have 
$$f(x)G_{N}(t-x)=\sum_{j=1}^{\infty}f(x_{j})\phi_{j}(x)G_{N}(t-x_{j}).$$
Choosing $t=x$, it follows that
$$f(x)G_{N}(0)=\sum_{j=1}^{\infty}f(x_{j})\phi_{j}(x)G_{N}(x-x_{j})=:s_{f,\infty}(x).$$ Therefore,
$$|E_{1}(x)|=|f(x)-f(x)G_{N}(0)+s_{f,\infty}(x)-S_{f,\infty}(x)|.$$
Noting that $g_{N}(0)=1$ and 
$$g_{N}(x)-G_{N}(x)=\frac{1}{\sqrt{2\pi}}\int_{\pi-\delta}^{\infty}\widehat{g_{N}}(w)e^{iwx}dw+\frac{1}{\sqrt{2\pi}}\int_{-\infty}^{\delta-\pi}\widehat{g_{N}}(w)e^{iwx}dw,$$
we have $\sup_{x\in[-1,1]}|E_{1}(x)|\le E_{1,N}$. 


\end{proof}

\end{document}
